% =============================================================================
\section{Mathematical and Information-Theoretic Formulation}\label{sec:math}
% =============================================================================
\subsection{Complex Representation, Magnitude Projection, and Information}
Let $Z\in\mathbb{C}^{C\times H\times W}$ denote the prepared multicoil complex input and let $Y\in\{0,1\}$ denote the binary slice label. The magnitude-only representation is the deterministic map
\begin{equation}
\mathcal{M}(Z)=|Z|.
\label{eq:magnitude-map}
\end{equation}
For any componentwise phase field $\Phi=(\phi_j)_j$,
\begin{equation}
\mathcal{M}\!\left(\mathrm{diag}(e^{i\phi_j})Z\right)=\mathcal{M}(Z).
\end{equation}
Thus magnitude identifies a large family of complex signals that differ only through local phase. Since $\mathcal{M}$ is deterministic, the data-processing inequality gives
\begin{equation}
I(Y;\mathcal{M}(Z))\le I(Y;Z),
\label{eq:dpi}
\end{equation}
without asserting that the inequality is strict for this dataset. Equation~\eqref{eq:dpi} is therefore a statement about representational capacity, not evidence that phase must improve classification.

The nuisance transformation explicitly targeted by the training procedure is much smaller than arbitrary componentwise phase modification. A common global rotation acts as
\begin{equation}
\rho_\alpha(Z)=e^{i\alpha}Z,\qquad \alpha\in[-\pi,\pi].
\label{eq:u1-action}
\end{equation}
The diagnostic label should be unchanged under this transformation, motivating a model whose internal features may be equivariant while its final output is invariant.

\subsection{Image and Fourier Domains as Unitary Coordinate Systems}
Let $\mathcal{F}$ denote the centered orthonormal two-dimensional Fourier transform applied independently to each coil. Under orthonormal normalization,
\begin{align}
\|\mathcal{F}Z\|_2 &= \|Z\|_2,\\
\langle \mathcal{F}Z,\mathcal{F}W\rangle &= \langle Z,W\rangle.
\label{eq:parseval}
\end{align}
The Fourier transform is therefore unitary and invertible. Before robust scaling, $Z\mapsto\mathcal{F}Z$ is a unitary coordinate transformation and does not remove information from the prepared complex tensor. After representation selection, however, each representation receives its own sample-level robust magnitude scaling scalar, so the actual model inputs are not related only by a unitary transform. The Fourier transform nevertheless alters the geometry presented to a finite convolutional network. A local convolution in image coordinates and a local convolution in Fourier coordinates impose different locality priors, even though the underlying coordinate transformation is invertible. The input-domain factor consequently tests an architecture--representation interaction rather than an information-presence/absence contrast.

The global phase action commutes with the Fourier transform,
\begin{equation}
\mathcal{F}(e^{i\alpha}Z)=e^{i\alpha}\mathcal{F}(Z),
\end{equation}
so global-rotation considerations apply in both domains.

\subsection{Complex-Linear and Widely-Linear Operators}
For $z=x+iy$ and $W=A+iB$, standard complex multiplication is equivalent to the constrained real-linear map
\begin{equation}
\begin{bmatrix}\Re(Wz)\\\Im(Wz)\end{bmatrix}
=
\begin{bmatrix}A&-B\\B&A\end{bmatrix}
\begin{bmatrix}x\\y\end{bmatrix}.
\label{eq:real-linear-complex}
\end{equation}
This block structure is the algebraic signature of complex linearity. Equivalently, a map $T$ is complex-linear when
\begin{equation}
T(iz)=iT(z).
\label{eq:c-linear}
\end{equation}
A standard complex convolution therefore occupies a strict subset of all real-linear maps on the concatenated Cartesian components.

A general widely-linear map is
\begin{equation}
T(z)=W_1z+W_2z^*.
\label{eq:widely-linear-form}
\end{equation}
The conjugate pathway relaxes the complex-linearity constraint. In finite dimensions, any real-linear map between complex vector spaces can be written in this form. The special case $W_2=0$ recovers complex linearity. Under a global phase rotation,
\begin{equation}
T(e^{i\alpha}z)=e^{i\alpha}W_1z+e^{-i\alpha}W_2z^*,
\end{equation}
so a nonzero conjugate term is not generally $U(1)$-equivariant. The standard-versus-widely-linear factor therefore trades stronger symmetry structure for greater real-linear flexibility.

\subsection{Phase-Compatible and Phase-Selective Nonlinearities}
A pointwise nonlinearity $g:\mathbb{C}\rightarrow\mathbb{C}$ is globally phase-equivariant when
\begin{equation}
g(e^{i\alpha}z)=e^{i\alpha}g(z).
\label{eq:activation-equivariance}
\end{equation}
Any gate of the form
\begin{equation}
g(z)=z\,\psi(|z|)
\end{equation}
with real-valued $\psi$ satisfies Eq.~\eqref{eq:activation-equivariance}. The four nonlinearities in the factorial description are
\begin{align}
\operatorname{modReLU}(z)&=z\frac{\operatorname{ReLU}(|z|+b)}{|z|+\epsilon},\\
\operatorname{MagSiLU}(z)&=z\,\sigma\!\left(a(|z|-b)\right),\\
\operatorname{CReLU}(z)&=\operatorname{ReLU}(\Re z)+i\,\operatorname{ReLU}(\Im z),\\
\operatorname{Cardioid}(z)&=\frac{1}{2}\left(1+\cos(\arg z)\right)z.
\end{align}
Here $b$ is a learned channel-wise bias in modReLU, while $a>0$ and $b$ are the learned slope and threshold of the magnitude-gated SiLU. The first two gates are real functions of magnitude and are globally phase-equivariant, including the epsilon stabilizer in modReLU, because all scalar factors in their implemented formulas depend only on $|z|$. CReLU applies independent real ReLUs to the two Cartesian components, whereas cardioid is explicitly selective with respect to the real-axis phase. The magnitude-gated SiLU is a study-specific complex adaptation of a SiLU-style gate, not standard real SiLU applied to a complex tensor.

\subsection{Normalization as Geometry Control}
The implemented complex RMS normalization uses per-channel power averaged over batch and spatial dimensions,
\begin{align}
p_c^{\mathrm{batch}}&=\mathbb{E}_{b,u,v}\!\left[\Re(z_{b,c,u,v})^2+\Im(z_{b,c,u,v})^2\right],\\
p_c^{\mathrm{run}}&\leftarrow(1-m)p_c^{\mathrm{run}}+m p_c^{\mathrm{batch}},
\end{align}
where the second line is the running-power update during training. Its output is
\begin{equation}
\widetilde z_{b,c,u,v}=\begin{cases}
z_{b,c,u,v}\exp(\ell_c)(p_c^{\mathrm{batch}}+\epsilon)^{-1/2},&\text{training},\\
z_{b,c,u,v}\exp(\ell_c)(p_c^{\mathrm{run}}+\epsilon)^{-1/2},&\text{evaluation}.
\end{cases}
\label{eq:rmsnorm-general}
\end{equation}
The learned scale is positive because it is parameterized as $\exp(\ell_c)$. Because the power statistic depends only on magnitude, a global rotation of $z_c$ leaves the scale unchanged, and the normalized output rotates with the input. This exact running-power complex formulation is a study-specific adaptation of RMS normalization rather than an implementation taken verbatim from the cited real-valued RMSNorm definition.

Complex BatchNorm instead treats $(\Re z,\Im z)$ as a two-dimensional random vector. After centering, a covariance matrix
\begin{align}
\boldsymbol{\mu}_c&=
\begin{bmatrix}\mathbb{E}_{b,u,v}[\Re z_{b,c,u,v}]\\
\mathbb{E}_{b,u,v}[\Im z_{b,c,u,v}]\end{bmatrix},\\
V_c&=\begin{bmatrix}V_{rr}&V_{ri}\\V_{ri}&V_{ii}\end{bmatrix},\\
\widehat{\mathbf z}_c&=(V_c+\epsilon I)^{-1/2}(\mathbf z_c-\boldsymbol{\mu}_c),\\
\mathbf y_c&=
\begin{bmatrix}\gamma_{rr}&\gamma_{ri}\\\gamma_{ri}&\gamma_{ii}\end{bmatrix}
\widehat{\mathbf z}_c+
\begin{bmatrix}\beta_r\\\beta_i\end{bmatrix},
\end{align}
where $\mathbf z_c=[\Re z_c,\Im z_c]^\mathsf{T}$ and the covariance entries are computed from the centered components. The batch means and covariance are updated as running statistics during training and replaced by their stored values at evaluation. Thus the implementation whitens the joint real-imaginary covariance and applies a learned symmetric affine map; it is not independent normalization of the two components. In general, this learned whitening-plus-affine map need not commute with every global rotation. The normalization factor therefore compares a phase-compatible scalar power normalization against a more flexible covariance-aware Cartesian normalization.

\subsection{Pooling and Local Phase Aggregation}
For local complex activations $z_j=r_je^{i\phi_j}$, average pooling obeys
\begin{equation}
\left|\frac{1}{n}\sum_jr_je^{i\phi_j}\right|^2
=\frac{1}{n^2}\left(\sum_jr_j^2+2\sum_{j<k}r_jr_k\cos(\phi_j-\phi_k)\right).
\label{eq:complex-average-coherence}
\end{equation}
The cross terms show explicitly that the pooled magnitude depends on local relative phase. Average pooling is complex-linear and globally phase-equivariant, but it creates a new complex value through vector averaging.

Magnitude-max and magnitude-median pooling follow a different mechanism. In each $2\times2$ window, max pooling selects the actual complex element $z_j$ maximizing $|z_j|^2$. Median pooling applies the median operation to the four $|z_j|^2$ values and returns the actual complex element at the selected index; for four values, the implementation selects the lower median. These implementation-defined selection indices are determined by magnitude and are invariant to global phase, so the selected complex sample then rotates with the input. The operators preserve one observed phase rather than averaging several phases. Average pooling, by contrast, averages the real and imaginary components and is therefore the complex arithmetic mean.

\subsection{Streams and Interaction Operators}
The complex-only model retains complex features until the final magnitude readout. A dual-stream model constructs a real magnitude branch in parallel with the complex branch. If $h_c(Z)$ denotes complex features and $h_r(|Z|)$ real features, the final classifier acts on a concatenation of the form
\begin{equation}
\left[\mathrm{GAP}\{h_r(|Z|)\},\;\mathrm{GAP}\{|h_c(Z)|\}\right].
\end{equation}
The dual architecture therefore tests complementarity between explicitly invariant magnitude features and features learned before magnitude projection.

For modulus gating, the complex-to-real pathway is
\begin{equation}
g=\sigma\!\left(K*|h_c|\right),\qquad h_r^{\mathrm{out}}=h_r\odot g,
\end{equation}
where $K$ is a real convolution. The gate is applied after every paired block, including the final block, when this interaction is selected; there is no reverse real-to-complex gate. The gate is insensitive to a common rotation of the complex features. It is a study-specific cross-stream construction motivated by multiplicative feature gating, not an established operator transferred unchanged from prior work.

The holographic module uses complex query, key, and value projections. For token indices $i,j$ and embedding dimension $d$, its score, weights, and value aggregation are
\begin{align}
\ell_{ij}&=\frac{\Re\{q_i k_j^H\}}{\sqrt d}-\gamma\frac{\left\|\,|q_i|-|k_j|\,\right\|_2^2}{d},\qquad \gamma=1,\\
a_{ij}&=\operatorname{softmax}_{j}(\ell_{ij}),\\
\widetilde v_i&=\sum_j a_{ij}v_j.
\end{align}
The aggregated complex values are passed through an output projection, added residually to the complex input, and then normalized:
\begin{equation}
 h_c^{\mathrm{out}}=\operatorname{ComplexRMSNorm2d}\!\left(h_c+\mathcal{O}(\widetilde V)\right).
\end{equation}
In the complex branch this module is placed after Block~3 and its pooling and before Block~4 when selected. For simultaneous rotation $q\mapsto e^{i\alpha}q$, $k\mapsto e^{i\alpha}k$, both terms in $\ell_{ij}$ are unchanged; if the projections preserve the rotation, the attention weights are invariant and the value aggregation is equivariant. The magnitude-distance penalty and its combination with the Hermitian score are study-specific adaptations, not claims inherited from the cited attention literature.

\subsection{Finite-Grid Factorial Summaries}
For a performance quantity $A$, let $\mathcal{G}$ denote the 480-model parameter-matched complex grid. The stream/interaction coordinate is one combined factor, with
\begin{equation}
\begin{aligned}
s(q)&\in\mathcal{S}_{\mathrm{si}},\\
\mathcal{S}_{\mathrm{si}}&=\{\text{complex-only/none},\ \text{complex-only/holographic},\\
&\quad\text{dual/none},\ \text{dual/modulus gate},\ \text{dual/holographic}\}.
\end{aligned}
\end{equation}
Thus no separate Cartesian product of stream and interaction levels is implied. The main-effect factors are input domain, pooling, normalization, convolution, activation, and $s$. Let $N_{f,l}$ be the number of configurations at level $l$ of factor $f$, and let $N_{fg,l,m}$ be the number at joint levels $(l,m)$. The real baseline has 1,685,890 trainable scalar parameters; complex widths are adjusted per configuration, and accepted configurations lie within 1\% of that scalar budget. We define
\begin{align}
\bar A&=\frac{1}{|\mathcal{G}|}\sum_{q\in\mathcal{G}}A_q,\\
\bar A_f(l)&=\frac{1}{N_{f,l}}\sum_{\substack{q\in\mathcal{G}\\f(q)=l}}A_q,\qquad
\delta_f(l)=\bar A_f(l)-\bar A,\\
\bar A_{fg}(l,m)&=\frac{1}{N_{fg,l,m}}\sum_{\substack{q\in\mathcal{G}\\f(q)=l,\,g(q)=m}}A_q,\\
\delta_{fg}(l,m)&=\bar A_{fg}(l,m)-\bar A-\delta_f(l)-\delta_g(m).
\label{eq:factor-effects}
\end{align}
These are finite-grid descriptive summaries. They characterize the observed architecture surface; they are not population-level causal effects or multiple independent hypothesis tests.

\begin{table*}[t]
\centering
\caption{Algebraic and analytic role of the six factorial dimensions.}
\label{tab:factor-algebra}
\scriptsize
\begin{tabularx}{\textwidth}{L{2.1cm}L{3.1cm}L{5.0cm}X}
\toprule
Factor & Levels & Algebraic distinction & Analytic question represented in the grid \\
\midrule
Input domain & image; Fourier & Centered orthonormal FFT followed by representation-specific robust scaling & Does the CNN locality/aggregation prior align differently with image and frequency coordinates? \\
Pooling & max; median; average & Order-statistic selection versus complex-linear vector averaging & Should downsampling preserve one observed complex activation or combine local phase through interference terms? \\
Normalization & RMS; complex BatchNorm & Scalar power normalization versus $2\times2$ Cartesian whitening and affine remapping & Is rotation-compatible amplitude normalization sufficient, or is covariance adaptation beneficial? \\
Convolution & standard; widely linear & $\mathbb{C}$-linear map versus general $\mathbb{R}$-linear map $W_1z+W_2z^*$ & Does conjugate flexibility help beyond the stronger complex-linearity constraint? \\
Activation & modReLU; magnitude-SiLU; CReLU; cardioid & Phase-equivariant gates versus Cartesian or phase-selective nonlinearities & How much global-rotation structure should be preserved through nonlinear feature extraction? \\
Stream/interaction & five implemented levels & Complex-only processing, invariant magnitude branch, modulus gate, or Hermitian/magnitude attention & Are magnitude and complex features complementary, and do explicit interactions improve fusion? \\
\bottomrule
\end{tabularx}
\end{table*}
